%% kapa: exam=1 round=35 year=2024 period=2 form=A subject=appraisal_law subject_ko=감정평가관계법규 seq=32 number=32
%% source: Q-Net 감정평가사 기출문제 (cst00309) artlSeq=5213229 file=raw/1008/5213229/2교시 문제지 원본.pdf
%% conversion: kordoc 4.24.51 -> tools/gen1.py (auto); review=proofread-sonnet; flags=none
\begin{question}{32}
부동산등기법령상 대표자나 관리인이 있는 법인 아닌 사단(이하 ‘비법인사단’)에 속하는 부동산의 등기에 관한 설명으로 옳은 것은?
\begin{choices}
  \item 등기에 관하여는 비법인사단의 대표자나 관리인을 등기권리자 또는 등기의무자로 한다.
  \item 비법인사단의 부동산등기용등록번호는 소재지 관할 등기소의 등기관이 부여한다.
  \item 권리에 관한 등기를 비법인사단의 명의로 할 때에는 그 대표자나 관리인의 성명, 주소를 등기사항으로 기록하지 않아도 된다.
  \item 비법인사단은 사용자등록을 하고 등기에 관하여 전자신청을 할 수 있다.
  \item 비법인사단의 대표자나 관리인이 등기를 신청한 경우 등기관은 등기를 마치면 그 대표자나 관리인에게 등기필정보를 통지한다.
\end{choices}
\end{question}
\answer{5}
